\documentclass[11pt]{article}
\usepackage[a4paper,margin=0.82in]{geometry}
\usepackage{microtype}
\usepackage[T1]{fontenc}
\usepackage{lmodern}
\usepackage{parskip}
\usepackage[hidelinks]{hyperref}
\pagestyle{empty}
\setlength{\emergencystretch}{2em}
\setlength{\parskip}{0.45em}

\begin{document}

{\raggedright
\textbf{Cristiano Capone}\\
Istituto Superiore di Sanit\`a, 00161 Rome, Italy\\
\href{mailto:cristiano0capone@gmail.com}{cristiano0capone@gmail.com}\par}

\vspace{0.6em}
{\raggedright The Editors, \textit{Brain}\par}

\vspace{0.8em}

Dear Editors,

We are pleased to submit our manuscript, ``Optimal stimulation sites are not the
most affected: personalised models of resting-state fMRI in Alzheimer's disease,''
by Cristiano Capone, Enza Cece, Andrea Ciardiello, Guido Gigante, Evaristo
Cisbani and Maurizio Mattia, for consideration in \textit{Brain}.

Neuromodulation is being actively pursued in Alzheimer's disease, and every such
intervention requires a decision that is currently made on anatomical grounds:
which site to stimulate. Targets have been selected where degeneration and
network disruption are greatest, and clinical results have been equivocal. Whether
``where the disease is worst'' is the right selection criterion has never been
tested causally, because doing so requires asking a counterfactual question of an
individual patient's brain.

We address this with personalised generative models. For each participant we fit
a subject-specific, cross-subject-identifiable model whose free-running dynamics
reproduce that individual's own resting-state activity, and we then intervene on
it \emph{in silico}. Diagnostic labels come from ADNI's own diagnosis query and
are independently corroborated by cognition (MMSE $28.7\pm1.4$ in controls versus
$22.7\pm2.8$ in patients, $p=1.6\times10^{-13}$), and the discriminant is
unrelated to age, education or sex. Restoring a virtual patient's connectivity
toward the control template reverts the disease signature, establishing that it is
correctable in principle, but the required correction is distributed: no
single-node edit reproduces it.

The clinically relevant result concerns where a focal drive should act. A
single-site drive at the node whose connectivity is most altered fails to revert
the classification at any amplitude, whereas a drive at the site selected by its
\emph{effect on the disease discriminant} reverts almost every patient. The two
site sets are essentially disjoint, sharing one parcel out of $29$ and $25$, and
none at all when the comparison is restricted to cortex. The reason is
mechanistic: the most-affected nodes are the most weakly connected, and therefore
poor actuators, while therapeutic leverage sits at well-connected cortical nodes
that can redistribute a corrective perturbation. Targets are heterogeneous across
patients, and a real-time closed-loop controller reaches comparable efficacy at
substantially lower dose using only causally available information.

We are explicit about scope. This is an in-silico study; no patient was
stimulated, and complete reversal requires amplitudes above the physiological
range. Its contribution is a principle for target selection and a computable,
per-patient prescription: if the effective target is systematically not the most
affected region, then anatomically guided targeting may be mis-specified, and
target choice should be model-informed and individualised. We believe this speaks
directly to how neuromodulation trials in Alzheimer's disease are designed, and
therefore to the readership of \textit{Brain}.

The manuscript is original, is not under consideration elsewhere, and all authors
have approved its submission. The data derive from the Alzheimer's Disease
Neuroimaging Initiative and are used under its Data Use Agreement. We declare no
competing interests, and would gladly suggest suitable referees on request. We
thank you for considering our work.

\vspace{1em}
{\raggedright
Sincerely,\\[0.3em]
Cristiano Capone, on behalf of all authors\par}

\end{document}
